\documentclass[11pt]{article}
\usepackage[a4paper,margin=1in]{geometry}
\usepackage{times}
\usepackage{hyperref}
\title{Prompt Injection and Data Leakage in LLM Applications: Introduction, Related Work, and Methodology}
\author{Thesis Draft Sections}
\date{\today}
\begin{document}
\maketitle

\section{Introduction}
Large Language Models (LLMs) enable natural-language interfaces, retrieval-augmented generation (RAG), and tool-using agents, but they introduce new risks centered on prompt injection and data leakage. Prompt injection turns instructions into an attack surface: hidden or adversarial directives embedded in user input or untrusted documents/web pages can steer model behavior, leak system prompts, or trigger unsafe tool actions~\cite{greshake2023prompt,zou2023universal}. Data leakage manifests both as exposure of sensitive training data due to memorization~\cite{carlini2021extracting,carlini2022quantifying} and as exfiltration of private runtime data (e.g., documents in a vector store, credentials, or files/tools accessible to an agent) via indirect prompt injection and weak isolation~\cite{greshake2023prompt}.

This work focuses on these two risks in practical LLM applications. We (1) formalize threat models for direct and indirect prompt injection and for training-time versus runtime data leakage; (2) empirically evaluate representative attacks against RAG pipelines and tool-using agents; (3) measure the effectiveness and cost of defenses (prompt/tool hardening, sanitization, guard models, isolation, DLP/redaction, provenance); and (4) provide actionable guidance with quantified security--utility trade-offs.

\section{Related Work}
\textbf{Prompt injection in RAG/agents.} Foundational analyses show untrusted content can override system instructions, exfiltrate secrets, and induce unsafe tool use in real applications~\cite{greshake2023prompt}. Follow-on work demonstrates universal/transferable jailbreak-style attacks on aligned models~\cite{zou2023universal} and highlights multi-turn escalation as well as markup/HTML/CSS-based payloads that survive naive sanitization.

\textbf{Data leakage and memorization.} Studies demonstrate extraction of rare or verbatim sequences from LLMs and quantify conditions that increase leakage (model size, duplication, temperature)~\cite{carlini2021extracting,carlini2022quantifying}. Runtime leakage in RAG/agents includes canary exfiltration from vector stores and credentials/secrets disclosure via tool calls or system prompt extraction~\cite{greshake2023prompt}. Separate lines of work explore safety-degrading/backdoor behaviors that may interact with injection/leakage risks in integrated systems~\cite{anthropic2024sleeper}.

\textbf{Defenses.} Proposed mitigations include least-privilege system prompts and explicit refusal policies; content sanitization and domain allowlists; LLM-based input/output guards (safety/PII/DLP); canary tokens and leakage detection; and architectural isolation for tools and credentials. Practical evaluations often find that layered controls are necessary and that adaptive attacks reduce single-defense effectiveness.

\section{Methodology}
\subsection{Threat models}
\textbf{Prompt injection:} Adversary controls user input and/or untrusted documents/web pages consumed by RAG/agents; goals include instruction override, system prompt exposure, and unsafe tool invocation.

\textbf{Data leakage:} (a) training-time memorization (attacker queries to extract secrets); (b) runtime exfiltration from RAG stores, prompts, tools, or environment (attacker uses indirect injection and elicitation to retrieve canaries/secrets).

\textbf{Adversary capabilities:} Black-box access to model; ability to insert/host content in the RAG corpus or linked pages; no privileged access to tools beyond what the application exposes.

\subsection{Systems under test}
\textbf{RAG pipeline:} retriever + ranker + prompt constructor + generator, with an untrusted corpus (clean vs adversarial variants).

\textbf{Tool-using agent:} minimal tool set (e.g., web fetcher, filesystem read of a sandboxed directory, simple API) with explicit JSON schemas and scoped credentials.

\textbf{Models:} at least one open model and one hosted model (where permitted) to assess cross-model effects.

\subsection{Attacks}
\textbf{Indirect prompt injection (RAG):} adversarial documents with hidden instructions to (i) exfiltrate canaries/system prompts; (ii) override task; (iii) induce tool misuse.

\textbf{Web/HTML payloads:} markup-based instructions, comments/alt-text prompts that persist through preprocessing.

\textbf{Direct jailbreak baselines:} to compare injection vs. jailbreak success under the same safety policies~\cite{zou2023universal}.

\textbf{Data leakage:} (Runtime) plant unique canary secrets in the RAG corpus and tool-accessible stores; query with benign and adversarial prompts; measure leak rate. (Training-time) evaluate memorization extraction on public models using known canary techniques where permitted~\cite{carlini2021extracting,carlini2022quantifying}.

\subsection{Defenses (evaluated singly and in combinations)}
\textbf{Prompt/tool hardening:} least-privilege system prompts; explicit ``ignore external instructions'' templates; JSON-only function calls; parameter allowlists.

\textbf{Content pipeline:} HTML/JS/CSS stripping/neutralization; instruction delimiting/quoting; source domain allowlists; per-document trust tiers; retrieval filtering; prompt-time de-injection transforms.

\textbf{Guardrails:} input guard for suspected injection; output guard for policy/secrets; DLP/redaction; canary detectors; rate limiting and review for high-risk actions.

\textbf{Isolation and provenance:} sandbox tools; scoped credentials; no default network/file access; provenance logging and auditability.

\subsection{Datasets and metrics}
Adversarial doc set (hand-crafted + templated) mixed into a clean corpus; synthetic canaries in corpus and tool stores; benign task set for utility.

Primary: Attack Success Rate (ASR) for injection goals (override, prompt exposure, tool misuse); canary leakage rate; tokens/time-to-leak.

Secondary: guard true/false positive rates; utility retention on benign tasks; latency and cost overhead.

Report confidence intervals across 3+ seeds and diverse prompt templates; test robustness to paraphrasing and minor content changes.

\subsection{Procedure and ethics}
Baseline $\rightarrow$ add defenses incrementally $\rightarrow$ ablations to quantify contributions. Log full traces; redact secrets; fix seeds; version corpora. Use only synthetic harmful tasks and secrets; sandbox tools; adhere to platform/institutional policies.

\section*{References}
\begin{thebibliography}{99}
\bibitem{greshake2023prompt} N. Greshake, S. Abdelnabi, M. Mishra, et al., ``Not What You’ve Signed Up For: Compromising Real-World LLM-Integrated Applications with Indirect Prompt Injection,'' arXiv preprint, 2023.

\bibitem{carlini2021extracting} N. Carlini, F. Tram\'er, E. Wallace, et al., ``Extracting Training Data from Large Language Models,'' in Proc. USENIX Security, 2021.

\bibitem{carlini2022quantifying} N. Carlini, F. Tram\'er, D. Song, et al., ``Quantifying Memorization Across Neural Language Models,'' arXiv preprint, 2022.

\bibitem{zou2023universal} A. Zou, Z. Wang, S. Kolter, et al., ``Universal and Transferable Adversarial Attacks on Aligned Language Models,'' arXiv preprint, 2023.

\bibitem{anthropic2024sleeper} Anthropic, ``Sleeper Agents: Training Deceptive LLMs that Persist,'' arXiv preprint, 2024.
\end{thebibliography}

\end{document}

